\section*{Section 6 Concept Quiz}
\addcontentsline{toc}{section}{Section 6 Concept Quiz}

Use this as a short retrieval check after finishing the chapter. On the website, each item should accept an answer before revealing the hint and worked explanation.

\begin{bgcheck}{epsilon-linear-q1}{expression}{epsilon/3}[Concept check]
For \(f(x)=3x+1\) at \(x=2\), give a working choice of \(\delta\) in terms of \(\varepsilon\).
\begin{bghint}
\(|f(x)-7|=3|x-2|\).
\end{bghint}
\begin{bgreveal}
Choose \(\delta=\varepsilon/3\).
\end{bgreveal}
\end{bgcheck}

\begin{bgcheck}{epsilon-linear-q2}{expression}{epsilon/5}[Concept check]
For \(f(x)=5x\) at \(x=1\), give a working \(\delta\).
\begin{bghint}
\(|5x-5|=5|x-1|\).
\end{bghint}
\begin{bgreveal}
Choose \(\delta=\varepsilon/5\).
\end{bgreveal}
\end{bgcheck}

\begin{bgcheck}{epsilon-quadratic-q1}{expression}{min(1,epsilon/5)}[Concept check]
For the standard proof that \(x^2\to4\) as \(x\to2\), enter the usual safe choice of \(\delta\).
\begin{bghint}
First require \(|x-2|<1\), which makes \(|x+2|<5\).
\end{bghint}
\begin{bgreveal}
Choose \(\delta=\min\{1,\varepsilon/5\}\).
\end{bgreveal}
\end{bgcheck}

\begin{bgcheck}{formal-reading-q1}{choice}{output}[Concept check]
In \(|f(x)-L|<\varepsilon\), does \(\varepsilon\) control input distance or output distance?
\begin{bghint}
Look at which quantities are being compared.
\end{bghint}
\begin{bgreveal}
It controls output distance.
\end{bgreveal}
\end{bgcheck}

\begin{bgsummary}[title={Quiz use on the website}]
This page should report completion by concept, not merely a total score. A wrong answer should link back to the exact lesson that teaches the missing method. The same questions may appear in randomized numerical variants, but the canonical explanation should remain stable.
\end{bgsummary}
